% chapters/18_rag.tex
% Part of: AI / LLM / Agents / Hardware Notes

\section{RAG \& Long-context Retrieval}

\subsection{Naive RAG pipeline}

Retrieval-Augmented Generation (RAG) grounds LLM responses in external knowledge
without retraining model weights. The canonical pipeline has two phases:

\textbf{Offline (indexing):}
\begin{enumerate}
  \item Load source documents (PDFs, web pages, databases) via document loaders.
  \item Chunk documents into retrievable units (see \S\ref{chunking}).
  \item Embed each chunk with a dense embedding model.
  \item Store embeddings + raw text in a vector store.
\end{enumerate}

\textbf{Online (query time):}
\begin{enumerate}
  \item Embed the user query with the same model.
  \item Run approximate nearest-neighbour (ANN) search to find top-$k$ chunks.
  \item Prepend retrieved chunks as context in the prompt.
  \item LLM generates an answer conditioned on context.
\end{enumerate}

\textbf{LangChain components} covering this pipeline:
\begin{itemize}
  \item \textbf{Document loaders:} Dropbox, Google Drive, OneDrive, YouTube, PubMed,
    URLs, Airtable, Figma, Notion, Pandas, MongoDB, Trello, and 100+ more.
  \item \textbf{Text splitters:} \texttt{RecursiveCharacterTextSplitter} (default),
    sentence splitters, token-aware splitters, semantic splitters.
  \item \textbf{Vector stores:} 50+ integrations including Chroma, Pinecone, Weaviate,
    pgvector, FAISS; represent chunks as dense embeddings for low-latency queries.
  \item \textbf{Retrievers:} abstract interface accepting a query string, returning a
    \texttt{Document} list; decouples retrieval logic from the store backend.
  \item \textbf{Memory:} full conversation history, summarisation-based, or $n$-most-recent
    exchange buffers to maintain multi-turn context.
\end{itemize}

\subsection{Embedding models \& vector DBs}

% TODO

\subsection{Chunking strategies}
\label{chunking}

Chunking strategy is one of the highest-leverage decisions in a RAG system: it
determines the granularity of what can be retrieved and how much context each chunk
carries.

\textbf{Static approaches (no LLM required):}
\begin{itemize}
  \item \textbf{Fixed-size / sliding window:} split every $n$ tokens with an $m$-token
    overlap. Simple and predictable but ignores document semantics.
  \item \textbf{Sentence / paragraph:} split on linguistic boundaries. Better cohesion
    but variable lengths can overflow the retriever's embedding dimension budget.
  \item \textbf{Recursive character splitting:} try progressively smaller separators
    (\verb|\n\n|, \verb|\n|, space, character) until chunks fit the target size.
    The LangChain default; good general-purpose baseline.
  \item \textbf{Semantic splitting:} embed consecutive sentences and split where
    cosine distance between adjacent embeddings exceeds a threshold. Higher quality,
    more expensive.
\end{itemize}

\textbf{Agentic chunking:}
An LLM agent determines chunk boundaries dynamically by reading the document and
reasoning about semantic coherence and logical structure — treating chunking as a
comprehension task rather than a rule-based process. Produces higher-quality, more
coherent chunks at significantly greater token cost. Best suited for high-value
corpora where retrieval quality justifies inference spend.

\subsection{Retrieval methods: dense, sparse, hybrid}

The retrieval step is the most critical quality lever in a RAG system. Three main
approaches exist:

\textbf{Dense retrieval:}
Documents and queries are embedded into dense vector representations; retrieval
returns the top-$k$ nearest neighbours by cosine or dot-product similarity.
\begin{itemize}
  \item Good at: semantic matching, paraphrases, meaning-based queries
    (``what causes inflation'' retrieves relevant content even without the word
    ``inflation'').
  \item Weakness: can miss exact keyword matches (names, IDs, version numbers,
    unique terms). A query for ``v2.3.1 release notes'' may retrieve semantically
    related content about other releases.
\end{itemize}

\textbf{Sparse retrieval (BM25 / TF-IDF):}
Keyword-based scoring; documents ranked by term overlap with the query, weighted by
term frequency and inverse document frequency.
\begin{itemize}
  \item Good at: exact matches, named entities, product codes, technical terms.
  \item Weakness: no semantic flexibility; ``automobile'' and ``car'' are treated
    as unrelated terms.
\end{itemize}

\textbf{Hybrid retrieval:}
Combines dense and sparse scores (typically via reciprocal rank fusion or a linear
combination). This is the recommended default for production RAG systems because:
\begin{itemize}
  \item Dense catches semantic meaning; sparse catches exact terms.
  \item Most real-world queries contain both semantic intent and specific
    keyword constraints.
  \item Adding sparse to a dense system is cheap (BM25 is fast and requires no GPU).
\end{itemize}

\textbf{RAG vs fine-tuning vs in-context learning --- when to use which:}
\begin{center}
\begin{tabular}{lll}
\toprule
\textbf{Use RAG when} & \textbf{Use fine-tuning when} & \textbf{Use in-context when} \\
\midrule
Facts change often & Style/tone must be consistent & Task fits in the prompt \\
Private docs involved & Task patterns repeat at scale & Fast iteration needed \\
Citations required & Output format is rigid & No training data available \\
No retraining budget & Latency must be minimised & \\
\bottomrule
\end{tabular}
\end{center}

\subsection{Reranking \& advanced RAG}

\textbf{Advanced RAG} wraps the naive pipeline with pre- and post-retrieval
optimisation steps:

\begin{itemize}
  \item \textbf{Query rewriting:} reformulate the user query before embedding retrieval.
    HyDE (Hypothetical Document Embeddings) generates a hypothetical answer and embeds
    \emph{that} instead; step-back prompting abstracts the query to a higher-level principle.
  \item \textbf{Reranking:} a cross-encoder model (e.g.\ Cohere Rerank, BGE-Reranker)
    re-scores the top-$k$ ANN candidates against the query with full attention, then
    passes only top-$m$ to the context window.
  \item \textbf{Contextual compression:} a post-retrieval LLM filters each chunk to
    only the sentences most relevant to the query, reducing noise in the prompt.
  \item \textbf{Multi-query retrieval:} generate $n$ paraphrases of the query, retrieve
    for each, deduplicate (by hash or semantic similarity), then rerank the merged set.
\end{itemize}

\textbf{Agentic RAG} replaces the fixed retrieve-then-generate pipeline with an agent
that can iterate:

\begin{itemize}
  \item \textbf{Adaptive retrieval:} the agent decides \emph{whether} to retrieve at
    all, which knowledge source to query, and whether a retrieved result is sufficient
    to answer --- queries traditional RAG cannot express.
  \item \textbf{Iterative refinement:} if retrieval is insufficient, the agent
    reformulates the query and retrieves again, possibly from a different source.
  \item \textbf{Multi-source fusion:} the agent queries heterogeneous backends (vector
    DB, SQL, web search, REST API) and synthesises a coherent answer.
  \item \textbf{Semantic caching:} frequently retrieved query--result pairs are cached
    at the agent level to avoid redundant embedding lookups across turns.
  \item \textbf{Output validation:} the agent can verify retrieved facts against a
    secondary source before including them in the final response.
\end{itemize}

\begin{intuitionbox}[title=Traditional RAG vs Agentic RAG]
  Traditional RAG: one retrieval step, static rules, single knowledge base, no
  validation, cannot iterate. Agentic RAG: adaptive strategy, multi-source, iterative,
  self-validating, full tool-use loop. The cost is latency and token spend; the benefit
  is accuracy on complex multi-step queries.
\end{intuitionbox}

\subsection{Long-context vs retrieval trade-offs}

% TODO

